\documentclass[a4paper,11pt]{article}
\usepackage[T1]{fontenc}
\usepackage[utf8]{inputenc}
\usepackage[italian]{babel}
\usepackage{amsmath,amssymb}
\usepackage{graphicx}
\usepackage{geometry}
\usepackage{xcolor}
\usepackage{listings}
\geometry{margin=2.5cm}

\lstdefinelanguage{MiniZinc}{
  morekeywords={int,set,of,array,var,constraint,include,solve,satisfy,output,show},
  sensitive=true,
  morecomment=[l]{\%},
  morestring=[b]"
}

\lstset{
  language=MiniZinc,
  basicstyle=\ttfamily\small,
  keywordstyle=\color{blue}\bfseries,
  commentstyle=\color{teal!70!black}\itshape,
  stringstyle=\color{orange!70!black},
  numbers=left,
  numberstyle=\tiny\color{gray},
  stepnumber=1,
  numbersep=8pt,
  showstringspaces=false,
  frame=single,
  breaklines=true
}

\title{SAT: Schur's Lemma}
\date{}

\begin{document}
\maketitle

\section*{Problema: Distribuzione di Biglie in Urne}
Porre $n$ biglie, numerate da $1$ ad $n$, in $3$ urne di modo che nessuna urna contenga terne di biglie $x, y, z$ tali che $x + y = z$


\subsection*{1.1 Modellazione} 
Una formulazione del problema in SAT è la seguente:

\paragraph{Variabili}

\[
X_{i,j}
\qquad
 i=1,\dots,n,
\quad
 j=1,2,3
\]

La variabile $X_{i,j}$ è vera se e solo se la biglia $i$ viene posta nell'urna $j$

\paragraph{Vincoli}

Ogni biglia deve stare in almeno una urna:
\[
\Phi_{\mathrm{atleast}}
=
\bigwedge_{i=1}^{n}
(
X_{i,1}\lor X_{i,2}\lor X_{i,3}
)
\]


Ogni biglia deve stare in al massimo una urna:
\[
\Phi_{\mathrm{atmost}}
=
\bigwedge_{i=1}^{n}
\quad
\bigwedge_{1\leq j_1<j_2\leq 3}
(
\lnot X_{i,j_1}\lor \lnot X_{i,j_2}
)
\]

I due vincoli precedenti dicono che ogni biglia sta in esattamente una delle tre urne

Consideriamo due biglie $i_1$ e $i_2$ tali che
\[
1\leq i_1<i_2\leq n,
\qquad
 i_1+i_2\leq n
\]

La terza biglia della terna è $i_1+i_2$

Per ogni urna $j=1,2,3$ dobbiamo impedire che le tre biglie stiano nella stessa urna. Quindi:
\[
\Phi_{\mathrm{schur}}
=
\bigwedge_{\substack{
1\leq i_1<i_2\leq n\\
i_1+i_2\leq n
}}
\quad
\bigwedge_{j=1}^{3}
(
\lnot X_{i_1,j}
\lor
\lnot X_{i_2,j}
\lor
\lnot X_{i_1+i_2,j}
)
\]

La formula complessiva è:
\[
\Phi_n
=
\Phi_{\mathrm{atleast}}
\land
\Phi_{\mathrm{atmost}}
\land
\Phi_{\mathrm{schur}}
\]

\newpage

\subsection*{1.2 Istanziazione per n = 5}
$\Phi_5$ si ottiene istanziando la formula generale con $n=5$

\paragraph{Variabili}

\[
X_{i,j}
\qquad
i=1,\dots,5,
\quad
j=1,2,3
\]

\paragraph{Vincoli}

Ogni biglia deve stare in almeno una urna:
\[
\Phi_{\mathrm{atleast}}
=
\bigwedge_{i=1}^5
(X_{i,1}\lor X_{i,2}\lor X_{i,3})
\]

Ogni biglia deve stare in al massimo una urna:
\[
\Phi_{\mathrm{atmost}}
=
\bigwedge_{i=1}^5
\bigl(
(\lnot X_{i,1}\lor \lnot X_{i,2})
\land
(\lnot X_{i,1}\lor \lnot X_{i,3})
\land
(\lnot X_{i,2}\lor \lnot X_{i,3})
\bigr)
\]


Per $n=5$ le terne proibite sono:
\[
(1,2,3),\quad (1,3,4),\quad (1,4,5),\quad (2,3,5)
\]

Quindi:
\[
\Phi_{\mathrm{schur}}
=
\bigwedge_{j=1}^{3}
\Bigl(
(\lnot X_{1,j}\lor \lnot X_{2,j}\lor \lnot X_{3,j})
\land
(\lnot X_{1,j}\lor \lnot X_{3,j}\lor \lnot X_{4,j})
\land
(\lnot X_{1,j}\lor \lnot X_{4,j}\lor \lnot X_{5,j})
\land
(\lnot X_{2,j}\lor \lnot X_{3,j}\lor \lnot X_{5,j})
\Bigr)
\]

La formula complessiva è:
\[
\Phi_5
=
\Phi_{\mathrm{atleast}}
\land
\Phi_{\mathrm{atmost}}
\land
\Phi_{\mathrm{schur}}
\]


\end{document}
